\chapter{De los átomos a las galaxias: las escalas del universo}\label{ch:g9:atoms-to-galaxies}

Hace nueve años este libro empezó con tus cinco sentidos y las patas de
una mariquita. Se cierra con un viaje que ningún sentido puede hacer
solo: cuarenta y una potencias de diez, del corazón de un
\omterm{def:g9:atoms-to-galaxies:atom}{átomo} al borde del universo observable --- contigo, con bastante
exactitud, en el medio. Un último capítulo, un último secreto que se
debía desde una vieja broma y una escalera sobre la que apoyarse para
toda la física que viene.

\section{Hacia abajo: dentro de la materia}

\begin{definition}[Los átomos]\label{def:g9:atoms-to-galaxies:atom}
Las \omterm{def:g7:states-of-matter:molecule}{moléculas} de tu séptimo año están a su vez construidas: cada
\omterm{def:g7:states-of-matter:molecule}{molécula} es un montaje de \emph{átomos}\index{átomo} --- el
alfabeto de la naturaleza, un centenar de clases, cuyas combinaciones
deletrean todas las sustancias. Una \omterm{def:g7:states-of-matter:molecule}{molécula} de agua: dos átomos
de hidrógeno y uno de oxígeno. Un átomo, a su vez, tiene arquitectura:
un \emph{núcleo}\index{núcleo} minúsculo, pesado y con carga positiva
en el centro y, muy por fuera de él, una nube de
\emph{electrones}\index{electrón} --- portadores de carga negativa casi
sin \omterm{def:g4:weight-and-mass:weight}{peso}. Los tamaños de los átomos rondan los $10^{-10}$ \unit{m}; sus
núcleos, los $10^{-15}$ --- cien mil veces más pequeños todavía.
\end{definition}

\begin{proposition}[La materia es casi toda vacío]\label{prop:g9:atoms-to-galaxies:empty}
Amplía un \omterm{def:g9:atoms-to-galaxies:atom}{átomo} hasta el tamaño de un gran estadio y su
\omterm{def:g9:atoms-to-galaxies:atom}{núcleo} será un guisante en el círculo central --- los
\omterm{def:g9:atoms-to-galaxies:atom}{electrones}, mosquitos en la grada más alta: todo lo que hay entre
medias está vacío. La mesa de mármol, el yunque de hierro, tu propia
mano: abrumadoramente vacíos, con su solidez convertida en una negativa
eléctrica --- las nubes de \omterm{def:g9:atoms-to-galaxies:atom}{electrones} de los \omterm{def:g9:atoms-to-galaxies:atom}{átomos} repeliéndose
unas a otras hasta crear la ilusión de lo lleno. Es una de las frases
más asombrosas de la física, y es la aritmética pura de $10^{-10}$
frente a $10^{-15}$.
\end{proposition}

\begin{example}[La vieja broma, pagada]\label{ex:g9:atoms-to-galaxies:electron}
Los capítulos de electricidad te debían una pequeña broma, y aquí está.
Las cargas que desfilan por todos los cables de este libro son
\emph{\omterm{def:g9:atoms-to-galaxies:atom}{electrones}} --- y son negativos, así que desfilan del
\omterm{def:g3:first-electric-circuit:battery}{borne} $-$ de la \omterm{def:g3:first-electric-circuit:battery}{pila} hacia el $+$: precisamente al
revés del sentido convencional que fijaron los físicos un siglo antes
de que nadie pudiera preguntarles a los desfilantes. La convención se
mantuvo --- todas las leyes de tus circuitos funcionan perfectamente con
ella ---, pero la pequeña ironía queda ahí: en todos los esquemas que
has dibujado, la multitud de verdad camina en sentido contrario.
\end{example}

\begin{example}[Y más abajo todavía]\label{ex:g9:atoms-to-galaxies:deeper}
El \omterm{def:g9:atoms-to-galaxies:atom}{núcleo} tiene piezas propias, y su historia --- cómo se unen,
por qué algunos \omterm{def:g9:atoms-to-galaxies:atom}{núcleos} se parten (el \omterm{def:g5:heat-and-insulation:heat}{calor} de la central) o se
fusionan (el del Sol) --- lleva a la física hasta los $10^{-15}$
\omterm{def:g2:measuring-length:units}{metros} y más allá, hasta las estructuras más pequeñas que han
sondeado los humanos. El volumen de bachillerato abre el
\omterm{def:g9:atoms-to-galaxies:atom}{núcleo} como es debido; los años de universidad van aún más lejos.
Los peldaños de abajo de la escalera se siguen construyendo.
\end{example}

\section{Hacia arriba: dentro del cielo}

\begin{example}[El ascenso]\label{ex:g9:atoms-to-galaxies:ascent}
Sube ahora, potencia de diez a potencia de diez, pasando por los
lugares conocidos de nueve años. Tú, alrededor de $10^{0}$ \unit{m}. El
patio del colegio, $10^{2}$. La Tierra, $10^{7}$ de lado a lado --- la
canica del modelo del campo de deportes. La distancia a la
\omterm{def:g2:moon-in-the-sky:moon}{Luna}, $4 \times 10^{8}$; la del Sol, $1.5 \times 10^{11}$ --- ocho
minutos \omterm{def:g1:light-and-shadows:source}{luz}. La anchura del \omterm{def:g4:solar-system:family}{Sistema Solar}, aproximadamente
$10^{13}$; la \omterm{def:g4:solar-system:star}{estrella} más cercana, $4 \times 10^{16}$ --- cuatro
\omterm{def:g8:speed-of-light:lightyear}{años luz}. La \omterm{ex:g5:stars-night-sky:milkyway}{Vía Láctea}, $10^{21}$ \omterm{def:g2:measuring-length:units}{metros} de
ciudad; la galaxia vecina, a $2 \times 10^{22}$; y el cielo más
profundo explorado, cerca de $10^{26}$ \omterm{def:g2:measuring-length:units}{metros} --- con la
\omterm{def:g1:light-and-shadows:source}{luz} de sus galaxias más vieja que la Tierra que hay debajo de tu
silla.
\end{example}

\begin{omfigure}
\begin{tikzpicture}[scale=0.85]
  \draw[-Stealth, very thick] (0,0.6) -- (11.4,0.6);
  \foreach \x/\p in {0.4/-15, 1.7/-10, 3.0/-5, 4.3/0, 5.6/5,
      6.9/10, 8.2/15, 9.5/20, 10.8/25} {
    \draw[thick] (\x,0.45) -- (\x,0.75);
    \node[below, font=\footnotesize] at (\x,0.4) {$10^{\p}$};
  }
  % labels on two alternating rows, with thin leaders
  \foreach \x/\h/\name in {0.4/0.95/núcleo, 1.7/1.6/átomo,
      3.0/0.95/hormiga, 4.3/1.6/tú, 5.6/0.95/{un país},
      6.9/1.6/{distancia al Sol}, 8.2/0.95/{estrellas cercanas},
      9.5/1.6/{la galaxia}, 10.8/0.95/{cielo profundo}} {
    \draw[omExo!80, thin] (\x,0.75) -- (\x,\h);
    \node[above, font=\footnotesize] at (\x,\h) {\name};
  }
\end{tikzpicture}

{\small La escalera de las escalas, en \omterm{def:g2:measuring-length:units}{metros}: cuarenta y una
potencias de diez del \omterm{def:g9:atoms-to-galaxies:atom}{núcleo} al cielo profundo --- con el lector de
pie casi exactamente en el peldaño del medio. (La hormiga nos perdona
el redondeo.)}
\end{omfigure}

\begin{method}[Pensar en órdenes de magnitud]\label{met:g9:atoms-to-galaxies:oom}
La primera herramienta del físico ante cualquier pregunta nueva:
\begin{enumerate}
  \item expresa la magnitud en notación científica;
  \item quédate solo con la potencia de diez --- el \emph{orden de
        magnitud}\index{orden de magnitud};
  \item compara escaleras, no cifras: una \omterm{def:g4:solar-system:star}{estrella} a $10^{16}$
        \omterm{def:g2:measuring-length:units}{metros} frente a un \omterm{def:g4:solar-system:planet}{planeta} a $10^{11}$ se
        diferencian en cinco peldaños --- cien mil veces, digan lo que
        digan las primeras cifras;
  \item solo entonces, si hace falta, afina las cifras.
\end{enumerate}
Nueve años de juzgar pasos --- ``¿tiene sentido esta respuesta?'' ---
maduran en esta costumbre. El primerísimo capítulo del volumen de
bachillerato construye sobre ella su cultura de la medida.
\end{method}

\begin{example}[Aritmética de escalera]\label{ex:g9:atoms-to-galaxies:arithmetic}
¿Cuántos \omterm{def:g9:atoms-to-galaxies:atom}{átomos} abarcan la anchura de un lápiz? Lápiz, $10^{-2}$
\unit{m}; \omterm{def:g9:atoms-to-galaxies:atom}{átomo}, $10^{-10}$: ocho peldaños --- cien millones de
\omterm{def:g9:atoms-to-galaxies:atom}{átomos}, hombro con hombro. ¿Cuántas Tierras hasta el Sol? $1.5
\times 10^{11}$ entre $1.3 \times 10^{7}$: unos $10^{4}$ --- diez mil
canicas del modelo del campo de deportes, puestas en fila. La escalera
responde en segundos lo que las cuentas crudas de kilómetros entierran
en ceros.
\end{example}

\section{La vista desde el peldaño del medio}

\begin{remark}[Lo que han construido nueve años]\label{rem:g9:atoms-to-galaxies:built}
Mira escalera abajo y cuenta lo que tienes. En $10^{-10}$: las
\omterm{def:g7:states-of-matter:molecule}{moléculas} y los \omterm{def:g9:atoms-to-galaxies:atom}{átomos} --- y con ellos los tres estados,
el \omterm{def:g5:heat-and-insulation:heat}{calor}, el relevo del \omterm{def:g3:sound-around-us:sound}{sonido} y los desfilantes de la
corriente. Alrededor de $10^{0}$: las \omterm{def:g2:pushes-and-pulls:force}{fuerzas} y la
\omterm{def:g5:everyday-energy:energy}{energía}, los rayos rectos de la \omterm{def:g1:light-and-shadows:source}{luz} y las leyes de los
circuitos --- la física de tamaño humano de nueve años de experimentos.
De $10^{7}$ a $10^{22}$: la Tierra que gira, las \omterm{def:g3:sun-days-seasons:year}{estaciones} de un
\omterm{def:g6:sun-earth-moon:axisorbit}{eje} inclinado, la \omterm{def:g2:moon-in-the-sky:moon}{Luna} que cae y una sola ley de
\omterm{def:g9:gravitation:gravitation}{gravitación} gobernando por igual el huerto y la galaxia. Ninguna
otra asignatura del colegio te entrega una sola escalera del interior
de la materia al borde de lo visible --- y todos los peldaños se
subieron con instrumentos, honradez y una aritmética que puedes
comprobar tú mismo.
\end{remark}

\begin{remark}[Lo que espera arriba]\label{rem:g9:atoms-to-galaxies:ahead}
Y los asuntos pendientes son lo mejor de todo. ¿Qué ley exacta ata la
\omterm{def:g2:pushes-and-pulls:force}{fuerza} al movimiento? ¿Qué \emph{es} la \omterm{def:g1:light-and-shadows:source}{luz}, para correr a
$3 \times 10^{8}$ \omterm{def:g2:measuring-length:units}{metros} por segundo y llevar colores dentro?
¿Qué se parte dentro del \omterm{def:g9:atoms-to-galaxies:atom}{núcleo} --- y qué forja la \omterm{def:g1:light-and-shadows:source}{luz} del Sol?
¿Por qué la cuenta de la \omterm{def:g5:everyday-energy:energy}{energía} útil va siempre cuesta abajo? Todas
esas preguntas se han planteado con honradez y se han aplazado con
honradez en estas páginas; todas tienen respuesta, y las respuestas
están entre las cosas más finas que conoce nuestra especie. Empiezan en
el volumen de bachillerato, dentro de un año. Trae la escalera --- y la
costumbre de comprobarlo todo.
\end{remark}

\section{Ejercicios}

\begin{exercise}[$\star$]\label{exo:g9:atoms-to-galaxies:1}
Ordena por tamaño: \omterm{def:g7:states-of-matter:molecule}{molécula}, \omterm{def:g9:atoms-to-galaxies:atom}{átomo},
\omterm{def:g9:atoms-to-galaxies:atom}{núcleo}, hormiga. Da la potencia de diez de cada uno en
\omterm{def:g2:measuring-length:units}{metros}.
\end{exercise}

\begin{exercise}[$\star$]\label{exo:g9:atoms-to-galaxies:2}
Describe la arquitectura del \omterm{def:g9:atoms-to-galaxies:atom}{átomo} --- qué hay en el centro, qué lo
rodea y los dos tamaños que dan lugar a la \omterm{prop:g5:mirrors-reflection:image}{imagen} del estadio.
\end{exercise}

\begin{exercise}[$\star$]\label{exo:g9:atoms-to-galaxies:3}
Paga la vieja broma: ¿quién desfila de verdad por un cable, en qué
sentido --- y por qué sobreviven todas las leyes de los circuitos a la
revelación?
\end{exercise}

\begin{exercise}[$\star$]\label{exo:g9:atoms-to-galaxies:4}
¿Qué es un \omterm{met:g9:atoms-to-galaxies:oom}{orden de magnitud}? Da el \omterm{met:g9:atoms-to-galaxies:oom}{orden de magnitud} de: tu
estatura; la longitud del aula; el diámetro de la Tierra.
\end{exercise}

\begin{exercise}[$\star$]\label{exo:g9:atoms-to-galaxies:5}
¿En cuántos peldaños de la escalera se diferencian: el
\omterm{def:g9:atoms-to-galaxies:atom}{átomo} y el \omterm{def:g9:atoms-to-galaxies:atom}{núcleo}? ¿Tú y la Tierra? ¿La distancia al
Sol y la de la \omterm{def:g4:solar-system:star}{estrella} más cercana?
\end{exercise}

\begin{exercise}[$\star$]\label{exo:g9:atoms-to-galaxies:6}
``La mesa es casi toda espacio vacío''. Defiende la frase con los dos
números atómicos --- y explica qué hace que la mesa se note
sólida de todas formas.
\end{exercise}

\begin{exercise}[$\star$]\label{exo:g9:atoms-to-galaxies:7}
Aritmética de escalera: ¿cuántos \omterm{def:g9:atoms-to-galaxies:atom}{átomos}, puestos en fila, abarcan
más o menos un milímetro?
\end{exercise}

\begin{exercise}[$\star$]\label{exo:g9:atoms-to-galaxies:8}
Sitúa en la escalera, con una potencia de diez cada uno: un
\omterm{def:g8:speed-of-light:lightyear}{año luz}; la \omterm{ex:g5:stars-night-sky:milkyway}{Vía Láctea}; el cielo más profundo
explorado.
\end{exercise}

\begin{exercise}[$\star\star$]\label{exo:g9:atoms-to-galaxies:9}
El modelo del estadio: si el \omterm{def:g9:atoms-to-galaxies:atom}{núcleo} es un guisante de
\qty{1}{cm}, ¿cuánto mide de ancho el \omterm{def:g9:atoms-to-galaxies:atom}{estadio-átomo}? (Razón de
$10^{5}$ --- pásalo a \omterm{def:g2:measuring-length:units}{metros} y júzgalo frente a un estadio de
verdad.)
\end{exercise}

\begin{exercise}[$\star\star$]\label{exo:g9:atoms-to-galaxies:10}
Una gota de agua contiene unas $10^{21}$ \omterm{def:g7:states-of-matter:molecule}{moléculas}. Con la
comparación del mar de gotas de tu capítulo de séptimo --- o con tu
propia estimación de las gotas de todos los océanos (unos $10^{21}$
\omterm{def:g6:mass-and-volume:volume}{litros}, con $10^{5}$ gotas cada uno) ---, pesa la vieja afirmación
de que una gota tiene más \omterm{def:g7:states-of-matter:molecule}{moléculas} que gotas tienen los mares.
\end{exercise}

\begin{exercise}[$\star\star$]\label{exo:g9:atoms-to-galaxies:11}
Tu cuerpo está cerca del medio de la escalera: ¿cuántos peldaños hay
más o menos hacia abajo hasta el \omterm{def:g9:atoms-to-galaxies:atom}{átomo}, y hacia arriba hasta el
cielo más profundo? ¿Qué dice esa simetría sobre el alcance de nueve
años de colegio?
\end{exercise}

\begin{exercise}[$\star\star\star$]\label{exo:g9:atoms-to-galaxies:12}
El último ejercicio del libro. Elige tres capítulos cualesquiera de tus
nueve años --- uno de los cursos 1 a 3, uno de los cursos 4 a 6 y uno
de los cursos 7 a 9 --- y escribe, para cada uno, la única frase que le
dirías a tu yo más joven cuando empezaba ese capítulo: qué enseñaba de
verdad, visto desde lo alto de la escalera. (No hay página de
soluciones para una carta a uno mismo --- pero escríbela; el curso que
viene te lo agradecerá.)
\end{exercise}

\section{Problema: el viaje por potencias de diez}

\begin{problem}[{Problema de fin de semana --- la clase rueda ``El
viaje por potencias de diez''; un solo zum, cuarenta y un peldaños; la
reunión de guion}]\label{pb:g9:atoms-to-galaxies:1}
La película de la clase: un zum continuo desde una manta de picnic
hacia fuera hasta el cielo profundo y después hacia dentro hasta el
\omterm{def:g9:atoms-to-galaxies:atom}{núcleo} --- una potencia de diez por segundo de película. Tú
escribes el guion científico.

\medskip\noindent\textbf{Parte I --- Hacia fuera.}
El zum arranca encuadrando un \omterm{def:g2:measuring-length:units}{metro} de manta: $10^{0}$.
\begin{enumerate}
  \item ¿En qué segundos contiene el encuadre por primera vez: el
        patio entero del colegio ($10^{2}$); la Tierra entera
        ($10^{7}$); la \omterm{def:g6:sun-earth-moon:axisorbit}{órbita} de la \omterm{def:g2:moon-in-the-sky:moon}{Luna} ($10^{9}$)?
  \item ¿Qué segundo encuadra la distancia al Sol --- y qué dice el
        narrador que es la edad de la \omterm{def:g1:light-and-shadows:source}{luz} del Sol, según el capítulo
        de la \omterm{def:g1:light-and-shadows:source}{luz}?
  \item ¿Entre qué segundos cruza el encuadre el gran vacío en el
        que el \omterm{def:g4:solar-system:family}{Sistema Solar} ya se ha acabado pero la
        \omterm{def:g4:solar-system:star}{estrella} más cercana todavía no ha entrado? ¿Qué debería
        enseñar honradamente la pantalla ahí?
  \item El zum termina en el segundo $26$: ¿qué llena el encuadre y
        cómo de vieja es la \omterm{def:g1:light-and-shadows:source}{luz} más antigua que hay en él,
        comparada con la Tierra?
\end{enumerate}

\medskip\noindent\textbf{Parte II --- Hacia dentro.}
El zum se invierte y se zambulle en una hoja que hay sobre la manta.
\begin{enumerate}[resume]
  \item ¿En qué segundos pasa la película por: una hormiga
        ($10^{-2}$, siendo generosos); una célula ($10^{-5}$); una
        \omterm{def:g7:states-of-matter:molecule}{molécula} ($10^{-9}$); un \omterm{def:g9:atoms-to-galaxies:atom}{átomo} ($10^{-10}$)?
  \item Entre la nube de \omterm{def:g9:atoms-to-galaxies:atom}{electrones} del \omterm{def:g9:atoms-to-galaxies:atom}{átomo} y su
        \omterm{def:g9:atoms-to-galaxies:atom}{núcleo}, la película tiene que hacer zum durante cinco
        segundos más a través de ¿qué? ¿Qué dice el narrador sobre
        el estadio?
  \item El último fotograma de la película, en $10^{-15}$: su sujeto
        --- y la frase de cierre honrada del narrador sobre los
        peldaños de más abajo.
  \item Suma la \omterm{def:g2:measuring-time:duration}{duración} de la película: ¿cuántos segundos hacia
        fuera y cuántos hacia dentro, a partir del \omterm{def:g2:measuring-length:units}{metro} de la
        manta?
\end{enumerate}

\medskip\noindent\textbf{Parte III --- La reunión de guion.}
\begin{enumerate}[resume]
  \item Un profesor objeta que la película dedica un segundo al paso
        de la casa a la calle y un segundo al paso de la galaxia al
        cúmulo de galaxias --- ``¡pasos disparatadamente
        desiguales!''. Defiende los segundos iguales con la lógica de
        la escalera.
  \item El narrador quiere una frase recurrente para las dos
        direcciones --- algo cierto en todos los peldaños sobre el
        vacío. Redáctala (el \omterm{def:g4:solar-system:family}{Sistema Solar} y el \omterm{def:g9:atoms-to-galaxies:atom}{átomo}
        tienen que caber los dos debajo).
  \item Los créditos tienen que nombrar los dos instrumentos que
        conquistaron las dos direcciones. Nómbralos, y los capítulos
        de este libro que presentaron sus principios.
  \item Las últimas palabras de la película le pertenecen al propio
        curso: escribe el epílogo de dos frases --- nueve años, una
        escalera y lo que empieza el año que viene.
\end{enumerate}
\end{problem}
